\newpage
\phantomsection
\label{prf:every-element-lies-below-the-supremum}

\begin{remark*}[Return]
\hyperref[prop:every-element-lies-below-the-supremum]{Return to Proposition}
\end{remark*}

\begin{theorem*}[Every Element Lies Below the Supremum]
If $s = \sup A$, then for every $x \in A$, $x \le s$.
\end{theorem*}

\begin{proof}
\textbf{Professional Standard Proof.}~\\
Let $A \subseteq \mathbb{R}$ be a non-empty set that is bounded above, and let $s = \sup A$. By the definition of the supremum, $s$ is the least upper bound of $A$. A necessary condition for $s$ to be the supremum is that it satisfies $\operatorname{UpperBound}(s, A)$. By the definition of an upper bound, it follows that $x \le s$ for every $x \in A$.
\end{proof}

\begin{proof}
\textbf{Detailed Learning Proof.}~\\

\textbf{Step 1.} Identify the dual nature of the supremum. For a value $s$ to be $\sup A$, it must simultaneously be an upper bound of $A$ and the smallest of all such upper bounds.

\textbf{Step 2.} The condition $\operatorname{UpperBound}(s, A)$ means that $s$ acts as a ceiling for the entire set.

\textbf{Step 3.} Translate the predicate into an inequality. By the formal definition of an upper bound, every element $x$ belonging to $A$ must be less than or equal to that bound. Therefore, $x \le s$ for all $x \in A$.
\end{proof}

\begin{remark*}[Proof structure]
This is a direct proof utilizing the "Upper Bound" component of the supremum's definition.
\end{remark*}

\begin{remark*}[Dependencies]~\\
\begin{itemize}
  \item \hyperref[def:upper-bound]{Definition of Upper Bound}
  \item \hyperref[def:supremum]{Definition of Supremum}
\end{itemize}
\end{remark*}